\medskip

Au début de l'année 2021, une colonie d'oiseaux comptait $40$ individus. L'observation conduit à modéliser l'évolution de la population par la suite $\left(u_n\right)$ définie pour tout entier naturel $n$ par:

\[\left\{\begin{array}{l c l}
u_0&=&40\\
u_{n+1}&=&0,008u_n\left(200 - u_n\right)
\end{array}\right.\]

où $u_n$ désigne le nombre d'individus au début de l'année $(2021+n)$.

\medskip


\begin{enumerate}
\item Donner une estimation, selon ce modèle, du nombre d'oiseaux dans la colonie au
début de l'année 2022.
\end{enumerate}

On considère la fonction $f$ définie sur l'intervalle [0~;~100] par $f(x) = 0,008x(200 - x)$.

\begin{enumerate}[resume]
\item Résoudre dans l'intervalle [0~;~100] l'équation $f(x) = x$.
\item
	\begin{enumerate}
		\item Démontrer que la fonction $f$ est croissante sur l'intervalle [0~;~100] et dresser son tableau de variations.
		\item En remarquant que, pour tout entier naturel $n$,\, $u_{n+1} = f\left(u_n\right)$ démontrer par récurrence que, pour tout entier naturel $n$ :

\[0 \leqslant u_n \leqslant u_{n+1} \leqslant 100.\]

		\item En déduire que la suite $\left(u_n\right)$ est convergente.
		\item Déterminer la limite $\ell$ de la suite $\left(u_n\right)$. Interpréter le résultat dans le contexte de l'exercice.
	\end{enumerate}
\item On considère l'algorithme suivant:

\begin{center}
\fbox{\begin{tabular}{l}
def seuil(p) :\\
\qquad n=0\\
\qquad u = 40\\
\qquad while u < p :\\
\quad \qquad n =n+1\\
\quad \qquad u = 0.008*u*(200-u)\\
\qquad return(n+2021)\\ 
\end{tabular}}
\end{center}

L'exécution de seuil(100) ne renvoie aucune valeur. Expliquer pourquoi à l'aide de la question 3.
\end{enumerate}

\bigskip
